\chapter{Introduzione}
% \addcontentsline{toc}{chapter}{Introduzione al corso}

\begin{redbar}
L’\textbf{interazione uomo-macchina (Human-Computer Interaction o HCI)} è una disciplina multidisciplinare che si occupa del design, dell'implementazione e della valutazione di sistemi interattivi destinati all’uso umano. L'obiettivo centrale è progettare sistemi che siano utili, usabili e adottati dagli utenti, concentrandosi sulle loro necessità e sui compiti che devono svolgere.
\end{redbar}

Il corso si propone di fornire agli studenti una solida comprensione delle teorie, dei modelli e delle regole fondamentali per la progettazione e la validazione di interfacce usabili. L'obiettivo principale è garantire l'acquisizione delle competenze necessarie per progettare sistemi interattivi efficaci, con un forte focus sull'utente finale. Inoltre, verrà introdotta la metodologia Agile, che consente una progettazione iterativa e flessibile, adattandosi ai cambiamenti e ai feedback continui.

In questo contesto, il corso si propone di fornire una comprensione approfondita del processo di progettazione incentrato sull'utente. L’idea è quella di integrare le conoscenze teoriche con l’applicazione pratica, permettendo agli studenti di sviluppare sistemi interattivi che rispondano ai bisogni reali delle persone. Il corso copre anche l’importanza dell’ascolto e della raccolta dei bisogni degli utenti, nonché la valutazione dei sistemi sviluppati con gli utenti stessi.

\section{Interfaccia}
Il concetto fondamentale di interfaccia, visto come il punto d'incontro tra due sistemi indipendenti, in particolare tra l'essere umano e la macchina, ha come obiettivo principale garantire un’efficace comunicazione e controllo della macchina da parte dell’uomo. Le interfacce includono elementi hardware, come i dispositivi fisici, e software, che comprendono i programmi e i sistemi operativi che permettono l'interazione tra l’utente e la macchina.

Le interfacce, quindi, non sono semplici strumenti di input e output, ma costituiscono un vero e proprio ponte tra i due attori: da una parte l’essere umano, con le sue capacità interpretative e creative, e dall'altra la macchina, che segue in modo deterministico le istruzioni ricevute. Questa distinzione tra creatività umana e determinismo meccanico è essenziale per comprendere l'importanza di un’interfaccia efficace, che deve essere in grado di interpretare i comandi umani e restituire risultati utili e comprensibili.

\subsection{Storia delle interfacce}
Un interessante esempio storico della nascita delle interfacce viene dall’abaco, uno dei primi dispositivi di calcolo, utilizzato da civiltà come quella greca e cinese. Qui, l’interfaccia consisteva nel movimento delle dita sugli elementi computazionali, che fornivano un output visibile attraverso la loro disposizione. Con il progredire della tecnologia, le interfacce si sono evolute fino a includere macchine più sofisticate come la Pascalina, costruita da Blaise Pascal nel 1642 per eseguire operazioni aritmetiche, e la macchina di Leibniz per moltiplicazioni e divisioni. Queste interfacce utilizzavano leve e altri meccanismi fisici per manipolare gli elementi computazionali interni.

Un passo importante nella storia delle interfacce è rappresentato dalla macchina analitica di Charles Babbage, considerata il precursore dei computer programmabili. La contessa Ada Lovelace, figlia del poeta Lord Byron, scrisse i primi programmi per questa macchina, gettando le basi per la programmazione moderna.

Le interfacce hanno continuato a evolversi nel corso del XIX secolo con l'introduzione di schede perforate da parte di Hermann Hollerith, utilizzate per il censimento americano del 1890. Queste schede permettevano di codificare informazioni che venivano lette tramite circuiti elettrici, segnando un importante passo avanti verso la digitalizzazione dell'informazione.

Con l'avvento del computer ENIAC, nel 1945, si fece un ulteriore salto tecnologico. L'ENIAC era in grado di eseguire complessi calcoli balistici, e il programma veniva definito collegando circuiti elettrici tramite cavi e interruttori. Questo segnò l'inizio dell’era moderna dei computer, dove l’interfaccia tra uomo e macchina divenne sempre più sofisticata.

Un punto di svolta nella storia delle interfacce si ebbe nel 1984, con la rivoluzione introdotta dalla Apple attraverso la GUI (Graphic User Interface), che utilizzava il paradigma WIMP: Windows, Icons, Mouse, Pointer. Questo nuovo approccio semplificava notevolmente l'interazione con i computer, rendendoli accessibili a un pubblico più vasto.

Negli anni successivi, le interfacce hanno continuato a evolversi, passando dalle classiche interfacce basate su tastiera e mouse a quelle post-WIMP, che includono sistemi di realtà virtuale, interfacce basate su gesti, riconoscimento vocale e persino interfacce tattili e cerebrali. Queste nuove tecnologie mirano a rendere l'interazione con le macchine sempre più naturale e intuitiva.

In conclusione, il passaggio dalle interfacce utente tradizionali all'esperienza utente (UX) rappresenta un cambiamento di paradigma importante: non si tratta più solo di progettare strumenti efficaci, ma di creare esperienze che rispondano ai bisogni e ai desideri degli utenti, rendendo la tecnologia sempre più integrata nella vita quotidiana.

\section{Progettazione di sistemi interattivi}
Progettare un buon sistema interattivo significa seguire un processo iterativo, centrato sugli utenti e sui loro bisogni. Questo processo include diversi passaggi: la raccolta delle necessità, la definizione di principi di design, la creazione rapida di prototipi e la valutazione del sistema. Ogni fase ha l’obiettivo di garantire che il sistema soddisfi gli obiettivi di usabilità, evitando soluzioni che non siano intuitive o che generino errori da parte degli utenti.

L'interazione uomo-macchina, infatti, coinvolge due entità che si influenzano reciprocamente nel tempo: l'uomo e il computer. Per progettare un sistema efficace, bisogna considerare non solo il sistema informatico, ma anche i modelli mentali degli utenti e i loro processi cognitivi.

L'HCI è un campo multidisciplinare, che attinge da diverse aree per garantire che i sistemi siano progettati in modo ottimale. Alcune delle principali discipline coinvolte sono la psicologia, le scienze cognitive e l’ergonomia, che forniscono informazioni sulle capacità fisiche e mentali degli utenti. La sociologia contribuisce a comprendere il contesto sociale in cui avviene l’interazione, mentre l’ingegneria informatica e la grafica aiutano nella costruzione dei sistemi e delle interfacce.

\subsection{Il modello di Norman}
Uno dei modelli più conosciuti nel campo dell'HCI è il modello di interazione di Norman. Questo modello descrive il processo che avviene quando un utente interagisce con un sistema. Il primo passo è la definizione di un obiettivo, che l'utente deve raggiungere attraverso una serie di azioni. Dopo aver eseguito le azioni, l'utente osserva lo stato del sistema per capire se l'obiettivo è stato raggiunto.

L'interazione può essere influenzata da due fattori principali: il \textit{gulf of execution} (il divario tra ciò che l'utente vuole fare e ciò che il sistema permette di fare) e il \textit{gulf of evaluation} (il divario tra ciò che il sistema presenta e ciò che l'utente comprende). La progettazione efficace riduce questi divari, facilitando l'interazione tra l'utente e il sistema.

Uno degli aspetti cruciali della progettazione è la gestione degli errori umani. Secondo Norman, gli errori non dovrebbero mai essere considerati come colpa dell'utente, ma come risultato di una cattiva progettazione. Gli errori umani possono essere suddivisi in due categorie principali: \textit{slip} (quando l’utente sa cosa fare ma esegue l’azione sbagliata) e \textit{mistake} (quando l’utente non comprende bene il sistema e formula un obiettivo errato).

Un buon sistema interattivo dovrebbe anticipare e minimizzare la possibilità di errori, rendendo più facile per l'utente correggere eventuali errori e comprendere lo stato del sistema.

\subsection{Progettazione centrata sull’utente}
La progettazione centrata sull’utente (\textit{User-Centered Design} o \textit{UCD}) pone le necessità, i desideri e le limitazioni degli utenti finali al centro del processo di design. Il coinvolgimento attivo degli utenti durante tutte le fasi del processo di sviluppo assicura che il sistema finale sia più facile da imparare, più efficiente da usare e che riduca il numero di errori umani. Tuttavia, una delle sfide principali è trovare gli utenti giusti e motivati da coinvolgere nel processo, e tradurre le loro esigenze in requisiti di design.

Il design partecipativo va oltre l'UCD, coinvolgendo direttamente gli utenti nel processo di progettazione attraverso la creazione di scenari, schizzi e prototipi. Questo approccio è particolarmente utile quando si progetta per utenti esperti o popolazioni mature, in quanto richiede una collaborazione attiva e costante da parte degli utenti coinvolti.

Il design thinking è un approccio incentrato sull'innovazione, che combina le esigenze degli utenti, le possibilità tecnologiche e i requisiti di successo del business. Il processo non lineare e iterativo del design thinking prevede cinque fasi: empatia, definizione dei bisogni, ideazione, prototipazione e test.

D'altro canto, la progettazione agile è un metodo di sviluppo evolutivo che si basa su cicli di rilascio rapidi e iterazioni continue. Questo approccio favorisce la creazione di prototipi a basso costo e richiede ispezioni rapide dell’usabilità, cercando di ridurre al minimo i tempi di sviluppo.

Il design dei servizi rappresenta un cambiamento contemporaneo verso la progettazione di esperienze complete, piuttosto che di prodotti isolati. Un esempio tipico è quello del passaggio dalla vendita di automobili come prodotti, all’offerta di servizi di mobilità come Uber. Questo tipo di design si basa su cinque principi fondamentali: centralità dell’utente, co-creazione, sequenzialità, evidenziazione e approccio olistico.

\subsection{Usabilità e processo di design centrato sull'uomo}
Infine, l’usabilità è un concetto centrale nell’HCI. Definita dall'ISO come 
\begin{center}
\textit{"la misura in cui un sistema può essere usato da utenti specifici per raggiungere determinati obiettivi in modo efficace, efficiente e soddisfacente"}    
\end{center}
l'usabilità include diverse dimensioni: utilità, facilità di apprendimento, memorizzabilità, efficacia, efficienza, visibilità dello stato del sistema, gestione degli errori e soddisfazione dell'utente.

Il processo di progettazione centrato sull'utente si sviluppa attraverso diverse fasi: il needfinding (ovvero la raccolta dei bisogni), l’analisi, il design, l'iterazione e prototipazione, e infine l’implementazione e il dispiegamento del sistema.